\section{Descriptive Statistics}
\label{sec:descriptive}

Figure~\ref{fig:cecl_equity_distribution} plots the cross-sectional
distribution of \texttt{CECL Adj}.  The distribution is right-skewed:
the full-sample mean is 0.28 percent of equity, the median is
essentially zero, and roughly half of community banks adopted CECL with
loan portfolios whose forward-looking expected credit losses required
little incremental reserve above the incurred-loss standard.  The right
tail is substantial: the High-CECL threshold at 2.5 percent of equity
sits at approximately the 95th percentile, and the 217 banks that
exceed it record a mean adjustment of 4 percent of equity.
Section~\ref{sec:robustness:scaling} examines the full shape of the
dose-response, including the left tail of negative adjustments.

Figure~\ref{fig:cecl_adoption_quarter_distribution} shows the
distribution of adoption quarters.  Adoption is concentrated in 2023Q1
(3,997 of 4,320 sample banks, 92.5 percent); the remaining 323 banks
adopt in adjacent quarters due to non-calendar fiscal years or minor
filing lags.  We use bank-specific event time throughout, so off-cycle
adopters contribute symmetric pre- and post-adoption windows.

Table~\ref{tab:descriptive_sumstats_cecl} compares High- and Low-CECL
banks at the adoption cross-section.  High-CECL banks enter with weaker
loan quality, higher NPL and allowance ratios, modestly higher CRE
concentration, and thinner capital buffers, consistent with more
credit-sensitive portfolios requiring larger expected-loss reserves.
Pre-adoption deposit structure is broadly similar on dimensions central
to our outcome variables.  Uninsured time deposits do not differ
significantly; the total uninsured deposit share is, if anything,
slightly \emph{lower} at high-CECL banks---the opposite of what a
run-exposure story would require.  High-CECL banks hold about 1.9
percentage points more insured time deposits at baseline, reflecting
greater reliance on retail certificate funding; bank fixed effects
absorb this level, and the pre-period event-study coefficients on both
series are flat.

Table~\ref{tab:cecl_predictors_cross_section} regresses \texttt{CECL Adj}
on 2022Q4 bank characteristics.  Column~(1) uses balance-sheet and
credit-quality fundamentals; column~(2) adds deposit composition and
eight-quarter pre-adoption performance trends; column~(3) adds local
deposit-market concentration.  The adjustment loads on NPL ratios, size,
capital, and deteriorating profitability trends in the expected
directions.  Among funding variables, the time-deposit and brokered
shares enter positively but with small coefficients, and the uninsured
deposit share---the variable most relevant to run exposure---is
unrelated to the adjustment.  Unrealized securities losses lose
significance once expanded controls enter.  Adjusted $R^2$ reaches 1.4
percent in the most saturated specification.  As discussed in
Section~\ref{sec:identification}, we do not treat the low $R^2$ as proof
of exogeneity; column~(2) serves as the first stage for the residualized
treatment used in Section~\ref{sec:robustness}.
